El presente proyecto tuvo como propósito mejorar y automatizar un brazo robótico cartesiano existente para su aplicación en la línea de empaque del Laboratorio de Diseño de Procesos de la Universidad del Valle de Guatemala. El sistema existente presentaba limitaciones mecánicas en los ejes X, Y y Z, tales como fricción, desalineaciones y movimientos irregulares, las cuales afectaban su precisión, repetibilidad y funcionamiento.

Para solucionar estas deficiencias, se realizó una reingeniería parcial de la estructura y de los sistemas de traslación, mediante la incorporación de ejes guía y rodamientos lineales. Además, se integraron un sistema de visión computacional, finales de carrera, un codificador rotativo y una arquitectura de control embebido para ejecutar tareas de detección, clasificación, manipulación y transferencia de productos.

El sistema también contó con un procedimiento de calibración automática, un modo de operación manual mediante un mando inalámbrico y una interfaz local para visualizar su estado. Finalmente, se evaluaron la precisión, repetibilidad, tiempo de respuesta y capacidad de detección del sistema mediante pruebas realizadas en un entorno académico controlado.